\documentclass[11pt]{article}

\usepackage{../ppackage}
\usepackage{bbm}
\usepackage{import}






\usepackage{tikz}
\usetikzlibrary{arrows.meta,patterns}
\usetikzlibrary{ipe} % ipe compatibility library

\usepackage{../../../Notes/tikzit}
\usepackage{../../../Notes/utf8math}

\input{../../../Notes/TIKZ/digraph.tikzstyles}



\usepackage{url}


\newcommand{\solution}{
\bigskip\noindent
	\textbf{Solution: \\}}
	


\DeclareMathOperator{\size}{size}
\DeclareMathOperator{\conv}{conv}
\newcommand{\SV}{\mathrm{SV}}
\newcommand{\bigO}{O}
\newcommand{\cut}{\mathrm{cut}}
\newcommand{\LLL}{\mathrm{LLL}}
\newcommand{\setR}{\mathbb{R}}
\newcommand{\setZ}{\mathbb{Z}}
\newcommand{\setQ}{\mathbb{Q}}
\newcommand{\setC}{\mathbb{C}}
\newcommand{\setN}{\mathbb{N}}
\newcommand{\wt}[1]{\widetilde{#1}}
\newcommand{\opt}{{\sc 0/1-opt}\xspace}
\newcommand{\aug}{{\sc 0/1-aug}\xspace}
\newcommand{\psep}{{\sc 0/1-psep}\xspace}
\newcommand{\sep}{{\sc 0/1-sep}\xspace}
\newcommand{\fopt}{{\sc 0/1-testopt\xspace} }

\newcommand{\hpp}{\mathrm{HPP}}
\newcommand{\nodes}{\mathcal{V}}
\newcommand{\vol}{\mathrm{vol}}
\newcommand{\diag}{\mathrm{diag}}
\newcommand{\arcs}{\mathcal{A}}
\newcommand{\edges}{\mathcal{E}}
\newcommand{\paths}{\mathscr{P}}
\newcommand{\cycles}{\mathcal{C}}




\newcommand{\K}{{\mathcal K}}
\newcommand{\A}{{A}}
\newcommand{\B}{{B}}
\newcommand{\T}{\mathscr{T}}
\newcommand{\eE}{\mathscr{E}}
\newcommand{\eS}{\mathscr{S}}
\newcommand{\eP}{\mathscr{P}}
\newcommand{\eM}{\mathscr{M}}



\newcommand{\transp}{^{\mathrm{T}}}

\newcommand{\smallmat}[1]{\left( \begin{smallmatrix} #1 \end{smallmatrix}\right)}

\newcommand{\mat}[1]{ \begin{pmatrix} #1 \end{pmatrix}}
\newcommand{\smat}[1]{ \big(\begin{smallmatrix} #1 \end{smallmatrix}\big)}

\newcommand{\pc}{\mathscr{P}}
\newcommand{\ob}{\mathscr{O}}
\newcommand{\odds}{\mathscr{W}}
\newcommand{\up}{\mathscr{U}}
\newcommand{\ef}{\mathscr{F}}
\newcommand{\eh}{\mathscr{H}}
\newcommand{\ev}{\mathscr{V}}
\newcommand{\ec}{\mathscr{C}}
\newcommand{\eu}{\mathscr{U}}

\newcommand{\lex}{\mathrm{lex}}

\renewcommand{\leq}{\leqslant}
\renewcommand{\geq}{\geqslant}









\newcommand{\linhull}{\mathrm{lin.hull}}
\newcommand{\affhull}{\mathrm{affine.hull}}
\newcommand{\charcone}{\mathrm{char.cone}}
\newcommand{\cone}{\mathrm{cone}}
\newcommand{\rank}{\mathrm{rank}}
\newcommand{\wb}[1]{\overline{#1}}



\usepackage{enumerate}

      
\institute{\'Ecole Polytechnique F\'ed\'erale de Lausanne}
\lecture{Discrete Optimization}
\faculty{Prof. Friedrich Eisenbrand}
\term{Spring 2026}
\publishdate{May 26, 2026}
\duedate{ }
\problemset{Problem Set~13 (Solutions)}

\begin{document}
\makeheader

\begin{enumerate}[1)]

\item \label{arbitrage} Let \(C\) be a set of currencies and let \(r_{ij}\) denote the exchange
rate from currency \(i\) to currency \(j\). We say that there is an
\emph{arbitrage opportunity} if there exists a sequence 
\[
i_0,i_1,\ldots,i_k=i_0 ∈ C
\]
such that
\[
\prod_{\ell=0}^{k-1} r_{i_\ell i_{\ell+1}}>1.
\]

 Explain how Bellman--Ford can be used
 to detect such an opportunity.

\begin{solution}
First construct a directed graph \(G\) whose vertices are the currencies in \(C\) and whose arcs are between any pair of currencies \(i,j\) (assuming \(r_{ij}>0\)). 
Assign each arc \(ij\) its weight \(c_{ij}=-\log(r_{ij})\). 
Then, there is an arbitrage opportunity if and only if there is a cycle in \(G\) whose total weight is negative,
since
\[
\left(\prod_{\ell=0}^{k-1} r_{i_\ell i_{\ell+1}}\right) > 1 \quad \Leftrightarrow \quad -\log\left(\prod_{\ell=0}^{k-1} r_{i_\ell i_{\ell+1}}\right) = \sum_{\ell=0}^{k-1} -\log(r_{i_\ell i_{\ell+1}}) = \sum_{\ell=0}^{k-1} c_{i_\ell i_{\ell+1}} < 0.
\]
This can be detected by the Bellman-Ford algorithm.
\end{solution}

\item \label{Floyd} (Floyd-Warshall algorithm) Let $G = (V,A)$ be a weighted directed graph with weights $c: A ⟶ ℝ$.
Define $D^{(k)}_{ij}$ to be the length of a shortest path from \(i\) to \(j\) whose \emph{internal}
vertices are only allowed to belong to $\{1,\dots,k\}$.
For \(k=0\), no internal vertices are allowed. Therefore
\[
D^{(0)}_{ij}
=
\begin{cases}
0, & i=j,\\
c_{ij}, & ij\in A,\\
+\infty, & \text{otherwise.}
\end{cases}
\]

\begin{enumerate}[i)] 
\item
 Suppose that $G$ does not contain a negative weight cycle. 
  Show the recurrence
  \[
    D^{(k)}_{ij}
    =
    \min\left\{
      D^{(k-1)}_{ij},
      D^{(k-1)}_{ik}+D^{(k-1)}_{kj}
    \right\}.
  \]
\item Prove that the matrix $D^{(n)}$ can be computed in time $O(n^3)$ steps.
\item  Show that $G$ contains a negative cycle if and only if there exists an $i ∈V$ such that $D^{(n)}_{ii} <0$. 
\end{enumerate}

\begin{solution}
\begin{enumerate}[i)]
\item
For any \(i,j\in V\), consider a shortest path \(P\) from \(i\) to \(j\) whose internal vertices belong only to
\(\{1,\dots,k\}\). We first show that $D^{(k)}_{ij}
    \geq \min\left\{
      D^{(k-1)}_{ij},
      D^{(k-1)}_{ik}+D^{(k-1)}_{kj}
    \right\}$. There are two cases.
\begin{itemize}
  \item If \(k\) is not an internal vertex of \(P\), then all internal vertices of \(P\) belong to
  \(\{1,\dots,k-1\}\), so the length of \(P\) is at least \(D^{(k-1)}_{ij}\).
  \item If \(k\) is an internal vertex of \(P\), since \(G\) has no negative weight cycle, we may assume that
  \(P\) is simple. This is because removing a cycle cannot increase the length of the path. Hence \(k\) appears exactly once on
  \(P\). We can then split \(P\) at \(k\) into an \(i\)-\(k\) path and a \(k\)-\(j\) path. The internal vertices of both
  subpaths belong only to \(\{1,\dots,k-1\}\). Therefore the length of \(P\) is at least $D^{(k-1)}_{ik}+D^{(k-1)}_{kj}.$
\end{itemize}
Next we show that $D^{(k)}_{ij}
    \leq \min\left\{
      D^{(k-1)}_{ij},
      D^{(k-1)}_{ik}+D^{(k-1)}_{kj}
    \right\}$.
Any path realizing \(D^{(k-1)}_{ij}\) is obviously a feasible path between \(i\) and \(j\) using only vertices from \(\{1,\dots,k\}\), and concatenating a path realizing
\(D^{(k-1)}_{ik}\) with a path realizing \(D^{(k-1)}_{kj}\) gives an \(i\)-\(j\) walk whose internal vertices are allowed to belong to \(\{1,\dots,k\}\). Since there are no negative cycles, this \(i\)-\(j\) walk can be turned into a \(i\)-\(j\) path of no
larger length. 

\item One can compute $D^{(n)}$ recursively. First we initialize the matrix \(D^{(0)}\) which takes \(O(n^2)\) time.
% This takes \(O(n^2)\) time if the matrix representation is used.
Then for the recursive step, assume that we have computed \(D^{(k-1)}\) for \(k\in\{1,\dots,n\}\). We can compute \(D^{(k)}\) from \(D^{(k-1)}\) using the recurrence in part i). This takes \(O(n^2)\) time, since there are \(n^2\) entries to update and each update takes \(O(1)\) time.
Let $T(k)$ be the time to compute \(D^{(k)}\). We have the recurrence
\[
T(k) = T(k-1) + O(n^2),
\]
which solves to \(T(n)=O(n^3)\).

\item
First if \(D^{(n)}_{ii}<0\) for some vertex \(i\), then the length of a shortest path from \(i\) to itself is negative, which gives us a negative cycle. 
Conversely, suppose that \(G\) contains a negative cycle \(C\). Choose a vertex \(i\) on \(C\). If we start at \(i\)
and go once around \(C\), we obtain a closed walk from \(i\) to \(i\) of negative total weight.
Therefore
\[
D^{(n)}_{ii}\leq c(C)<0.
\]
Hence \(D^{(n)}_{ii}<0\).
\end{enumerate}
\end{solution}

\item \label{hamilton} (Hamilton path problem) The following problem is NP-complete. So far, no polynomial-time algorithm has been found to solve it:

  Given a directed graph $G = (V,A)$ and two vertices $s≠t ∈ V$, decide whether there exists a path from $s$ to $t$ of length $|V|-1$. This is a path that spans all the vertices of $G$.

  \begin{enumerate}[i)] 
  \item  Show that this problem can be solved in polynomial time, if the graph does not have any directed cycles.
  \item Conclude that the shortest path problem (under presence of negative  cycles) should be difficult. 
  \end{enumerate}

\begin{solution}
\begin{enumerate}[i)]
\item 
For a directed acyclic graph \(G\), define the \emph{topological ordering} of $V$ to be a linear ordering of the vertices such that for every arc \(uv\) in \(G\), \(u\) appears before \(v\) in the ordering.

\textbf{Algorithm}.
\begin{enumerate}[1.]
  \item Compute a topological ordering of the vertices of \(G\). This can be done as follows (known as Kahn's algorithm):
  \begin{itemize}
    \item Initialize a queue \(Q\) with all vertices of in-degree zero.
    \item Initialize an empty list \(L\) to store the topological ordering.
    \item While \(Q\) is not empty:
    \begin{itemize}
      \item Remove a vertex \(u\) from \(Q\) and append it to \(L\).
      \item For each vertex \(v\) such that there is an arc \(uv\) in \(G\):
      \begin{itemize}
        \item Remove the arc \(uv\) from \(G\).
        \item If \(v\) has no other incoming arcs, add \(v\) to \(Q\).
      \end{itemize}
    \end{itemize}
  \end{itemize}
  \item Check whether \(s\) appears before \(t\) in the topological ordering. If not, then there is no Hamilton path from \(s\) to \(t\).
  \item If \(s\) appears before \(t\) in the topological ordering, then check whether there is an arc from each vertex to the next vertex in the topological ordering. If not, then there is no Hamilton path from \(s\) to \(t\). If yes, then there is a path from \(s\) to \(t\) of length \(|V|-1\).
\end{enumerate}

\textbf{Correctness}. 

We claim that $G$ has a Hamilton path from \(s\) to \(t\) if and only if \(s\) appears before \(t\) in the topological ordering and there is an arc from each vertex to the next vertex in the topological ordering.

Let's prove the "if" direction first. Assume that \(s\) appears before \(t\) in the topological ordering and there is an arc from each vertex to the next vertex in the topological ordering. Then we can construct a path from \(s\) to \(t\) by following the arcs in the order of the topological ordering. Since there are \(|V|\) vertices and we are following a path that visits each vertex exactly once, this path has length \(|V|-1\).
Next we prove the "only if" direction. Assume that there is a Hamilton path from \(s\) to \(t\). Since \(G\) is acyclic, the vertices on this path must appear in the order they are visited. Therefore, \(s\) must appear before \(t\) in any topological ordering of \(G\). Moreover, since the path visits each vertex exactly once, there must be an arc from each vertex to the next vertex in the topological ordering.

\textbf{Running time}. 

Computing a topological ordering takes \(O(|V|+|A|)\) time, and checking the conditions takes \(O(|V|)\) time. Therefore, the overall running time is \(O(|V|+|A|)\).  

\item Suppose not. Then we can solve the Hamilton path problem in polynomial time by solving the shortest path problem with negative cycles. Given an instance of the Hamilton path problem, we can assign a weight of \(-1\) to each arc and then check whether there is a path from \(s\) to \(t\) of length \(|V|-1\) by checking whether the shortest path from \(s\) to \(t\) has length \(-(|V|-1)\). This would contradict the NP-completeness of the Hamilton path problem. 
\end{enumerate}
\end{solution}

\item Consider the following problem. We are given $B ∈\mathbb{N}$, and a set of integer points $S= \{p ∈\setZ^n :
0 ≤p_i ≤B,\; ∀ i = 1,...,n\}$, whose points are all colored blue but one, which is red. We have an
oracle that, given vectors $l,r ∈\setR^n$, tells us whether the red point in $S$ is contained in the box
$S∩\{x ∈\setR^n : l_i ≤x_i ≤r_i,\;∀i = 1,...,n\}$ or not. Give an algorithm to find the red point using
$O(n \log(B))$ many oracle calls.


\begin{solution}
We split the problem in $n$ subproblems: for $i∈\{1,...,n\}$, we want to obtain the i-th component
of the red point. This is a simple binary search problem, and we illustrate how to solve it for $i= 1$:
we first call the oracle with $l = (0,...,0),r = (B/2,B,...,B)$. Then we call the oracle with $l=
(0,...,0),r = (B/4,B,...,B)$ (if the answer is positive) or $l = (B/2,...,0),r = (3B/4,B,...,B)$
(if the answer is negative), and so on. In this way, we are guaranteed to find each component of the
red point with $O(\log(B))$ oracle calls, for a total of $O(n\log(B))$ many oracle calls.
\end{solution}

\end{enumerate}




  

\end{document}

%%% Local Variables:
%%% mode: latex
%%% TeX-master: t
%%% End:
